
\exo

Parmi les fonctions définies par les expressions suivantes, reconnaître celles qui sont affines et préciser les valeurs de $m$ et $p$ :

\vspace{-8pt}
\setlength{\columnsep}{1cm}
\setlength{\columnseprule}{0pt}
\begin{multicols}{3}
\begin{enumerate}
\item $f(x)=-2x+1$
\smallskip
\item $f(x)=\dfrac{2x}{3}$
\item $f(x)=(2+x)(2x-1)$
\smallskip
\item $f(x)=\dfrac{1-2x}{3}$
\item $f(x)=x-(2x+1)$
\smallskip
\item $f(x)=\dfrac{2}{3x}$
\end{enumerate}
\end{multicols}

\exo

\begin{enumerate}
\item Écrire l'expression de la fonction $f$ qui, à un prix $x$, associe le prix $f(x)$ après une hausse de \SI[retain-explicit-plus]{20}{\percent}.
\item Reprendre cette question dans le cas où la hausse est de $t\;\%$.
\end{enumerate}

\exo

Représenter, dans un même repère, les fonctions affines suivantes :
\[
f:x\mapsto -2x+3
\qpvq
g:x\mapsto \frac{1}{2}x-4
\qpvq
h:x\mapsto 2-x
\qpvq
k:x\mapsto 3x-3.\]

\exo

Associer chaque droite à la fonction qu'elle représente :
\begin{center}
\def\xmin{-4}
\def\xmax{4}
\def\ymin{-4}
\def\ymax{4}
\def\xunite{0.5cm}
\def\yunite{0.5cm}
\psset{unit=\xunite, xunit=\xunite, yunit=\yunite, algebraic=true, dimen=middle, dotstyle=*, dotsize=3pt 0, linewidth=0.8pt, arrowsize=3pt 2, arrowinset=0.25, comma=true}
\begin{pspicture*}(\xmin,\ymin)(\xmax,\ymax)
\psgrid[xunit=\xunite, yunit=\yunite, subgriddiv=0, gridlabels=0, gridcolor=lightgray](0,0)(\xmin,\ymin)(\xmax,\ymax)
\psaxes[labels=none,labelFontSize=\scriptstyle, xAxis=true, yAxis=true, Dx=1, Dy=1, ticksize=-2pt 0, subticks=1]{->}(0,0)(\xmin,\ymin)(\xmax,\ymax)
\psplot[linewidth=1pt, plotpoints=200]{\xmin}{\xmax}{3*x-2}
\psdots(0,-2)(1,1)(2,4)
\begin{footnotesize}
\uput{3pt}[090](-3,5.75){$d_1$}
\uput{5pt}[225](0,0){$0$}
\uput{4pt}[270](1,0){$1$}
\uput{4pt}[180](0,1){$1$}
\end{footnotesize}
\end{pspicture*}
\quad
\begin{pspicture*}(\xmin,\ymin)(\xmax,\ymax)
\psgrid[xunit=\xunite, yunit=\yunite, subgriddiv=0, gridlabels=0, gridcolor=lightgray](0,0)(\xmin,\ymin)(\xmax,\ymax)
\psaxes[labels=none,labelFontSize=\scriptstyle, xAxis=true, yAxis=true, Dx=1, Dy=1, ticksize=-2pt 0, subticks=1]{->}(0,0)(\xmin,\ymin)(\xmax,\ymax)
\psplot[linewidth=1pt, plotpoints=200]{\xmin}{\xmax}{-2*x+3}
\psdots(0,3)(1,1)(2,-1)(3,-3)
\begin{footnotesize}
\uput{3pt}[090](-3,5.75){$d_1$}
\uput{5pt}[225](0,0){$0$}
\uput{4pt}[270](1,0){$1$}
\uput{4pt}[180](0,1){$1$}
\end{footnotesize}
\end{pspicture*}
\quad
\begin{pspicture*}(\xmin,\ymin)(\xmax,\ymax)
\psgrid[xunit=\xunite, yunit=\yunite, subgriddiv=0, gridlabels=0, gridcolor=lightgray](0,0)(\xmin,\ymin)(\xmax,\ymax)
\psaxes[labels=none,labelFontSize=\scriptstyle, xAxis=true, yAxis=true, Dx=1, Dy=1, ticksize=-2pt 0, subticks=1]{->}(0,0)(\xmin,\ymin)(\xmax,\ymax)
\psplot[linewidth=1pt, plotpoints=200]{\xmin}{\xmax}{0.667-0.333*x}
\psdots(-1,1)(2,0)(-4,2)
\begin{footnotesize}
\uput{3pt}[090](-3,5.75){$d_1$}
\uput{5pt}[225](0,0){$0$}
\uput{4pt}[270](1,0){$1$}
\uput{4pt}[180](0,1){$1$}
\end{footnotesize}
\end{pspicture*}
\end{center}
\vspace{-6pt}
\[
f(x)= 3-2x
\qpvq
g(x)= 3x-2
\qpvq
h(x)=-\frac{1}{3}x+\frac{2}{3}.
\qpvq
k(x)=\frac{2}{3}x-\frac{1}{3}.
\]

\exo

En observant le graphique ci-dessous, compléter le tableau :
\vspace{-6pt}

\begin{center}
\begin{minipage}[t]{7.9cm}
~

\vspace{-6pt}
\def\xmin{-4}
\def\xmax{9}
\def\ymin{-6}
\def\ymax{7}
\def\xunite{0.6cm}
\def\yunite{0.6cm}
\psset{unit=\xunite, xunit=\xunite, yunit=\yunite, algebraic=true, dimen=middle, dotstyle=*, dotsize=3pt 0, linewidth=0.8pt, arrowsize=3pt 2, arrowinset=0.25, comma=true}
\begin{pspicture*}(\xmin,\ymin)(\xmax,\ymax)
\psgrid[xunit=\xunite, yunit=\yunite, subgriddiv=0, gridlabels=0, gridcolor=lightgray](0,0)(\xmin,\ymin)(\xmax,\ymax)
\psaxes[labels=none,labelFontSize=\scriptstyle, xAxis=true, yAxis=true, Dx=1, Dy=1, ticksize=-2pt 0, subticks=1]{->}(0,0)(\xmin,\ymin)(\xmax,\ymax)[{\footnotesize $x$},135][{\footnotesize $y$},-45]
\psplot[linewidth=1pt, plotpoints=200]{\xmin}{\xmax}{5-0.25*x}
\psplot[linewidth=1pt, plotpoints=200]{\xmin}{\xmax}{-3+2*x}
\psplot[linewidth=1pt, plotpoints=200]{\xmin}{\xmax}{0.75*x}
\psplot[linewidth=1pt, plotpoints=200]{\xmin}{\xmax}{1-x}
\psplot[linewidth=1pt, plotpoints=200]{\xmin}{\xmax}{-3-0.25*x}
\psplot[linewidth=1pt, plotpoints=200]{\xmin}{\xmax}{3}
\begin{footnotesize}
\uput{3pt}[090](-3,5.75){$d_1$}
\uput{3pt}[180](-1,-5){$d_2$}
\uput{3pt}[135](8,6){$d_3$}
\uput{3pt}[045](-3,4){$d_4$}
\uput{3pt}[080](8,-5){$d_5$}
\uput{3pt}[270](-3.5,3){$d_6$}
\uput{5pt}[225](0,0){$0$}
\uput{4pt}[270](1,0){$1$}
\uput{4pt}[180](0,1){$1$}
\end{footnotesize}
\end{pspicture*}
\end{minipage}
\quad
\begin{minipage}[t]{8.9cm}
~

\vspace{-6pt}
\setlength{\tabcolsep}{0.1cm}
\renewcommand{\arraystretch}{1.5}
\begin{tabular}{
|>{\centering\arraybackslash}m{1.5cm}
|>{\centering\arraybackslash}m{2.0cm}
|>{\centering\arraybackslash}m{2.0cm}
|>{\centering\arraybackslash}m{3.0cm}|}
\hline 
\textbf{Droite} & \textbf{Coefficient directeur} & \textbf{Ordonnée à l'origine} & \textbf{Fonction associée} \\ 
\hline
$d_1$ &  &  &  \\
\hline
$d_2$ & $2$ & $-3$ & $x\mapsto 2x-3$ \\
\hline 
$d_3$ &  &  &  \\
\hline 
$d_4$ &  &  &  \\
\hline 
$d_5$ &  &  &  \\
\hline 
$d_6$ &  &  &  \\
\hline 
\end{tabular}
\end{minipage}
\end{center}

\exo

Vrai ou faux ? Justifier la réponse choisie.
\begin{enumerate}
\item Toute fonction linéaire est affine.
\item Toute fonction affine est linéaire.
\item Toute fonction constante est affine.
\item Toute fonction affine est représentée dans un repère par une droite.
\item Dans un repère, toute droite est la représentation d'une fonction affine.
\item Si une fonction $f$ est telle que $f(2)=3$ et $f(4)=6$, alors $f$ est linéaire.
\end{enumerate}

\exo

Dans chaque cas, déterminer la fonction affine $f$ à partir des informations données :

\vspace{-8pt}
\setlength{\columnsep}{1cm}
\setlength{\columnseprule}{0pt}
\begin{multicols}{2}
\begin{enumerate}
\item $f(0)=2$ et $f(2)=4$
\smallskip
\item $f(3)=2$ et $f(0)=-1$
\item $f(-2)=1$ et $f(6)=3$
\smallskip
\item $f(-1)=4$ et $f(1)=2$
\end{enumerate}
\end{multicols}
\vspace{-6pt}

\exo

Déterminer le sens de variation des fonctions affines définies sur $\R$ par :

\vspace{-8pt}
\setlength{\columnsep}{1cm}
\setlength{\columnseprule}{0pt}
\begin{multicols}{3}
\begin{enumerate}
\item $f(x)=2x+3$
\smallskip
\item $f(x)=-4x+5\vphantom{\dfrac{2}{7}}$
\item $f(x)=x+7$
\smallskip
\item $f(x)=8-x\vphantom{\dfrac{2}{7}}$
\item $f(x)=\sqrt{3}\;(x-2)$
\smallskip
\item $f(x)=\dfrac{3-2x}{7}$
\end{enumerate}
\end{multicols}
\vspace{-6pt}

\exo

\begin{enumerate}
\item La fonction affine $f$ est telle que $f(2)=5$ et $f(6)=3$.

La fonction $f$ est-elle strictement croissante ou strictement décroissante sur $\R$ ?
\item Reprendre la question précédente avec la fonction $g$ telle que $g(-1)=3$ et $g(2)=6$.
\end{enumerate}

\exo

On considère une fonction affine $f$, et on note $d$ sa droite représentative.

Parmi les situations suivantes, lesquelles sont impossibles ?
\begin{enumerate}
\item $f$ est strictement croissante sur $\R$, $d$ a pour ordonnée à l'origine $3$ et $f(2)=1$.
\item $f$ est strictement décroissante sur $\R$, $d$ a pour coefficient directeur $1$ et $f(2)=0$.
\item $f$ est strictement croissante sur $\R$, $d$ a pour ordonnée à l'origine $12$ et $f(7)=13$.
\item $f$ est strictement décroissante sur $\R$, $d$ a pour ordonnée à l'origine $0$ et $f(-3)=-3$.
\end{enumerate}

\exo

On considère une fonction affine $f$ pour laquelle on dispose du tableau de valeurs suivant :
\begin{center}
\setlength{\tabcolsep}{0.1cm}
\renewcommand{\arraystretch}{1.5}
\begin{tabular}{
|>{\centering\arraybackslash}m{2cm}
*{7}{|>{\centering\arraybackslash}m{1.5cm}}|}
\hline 
$x$ & & $-1$ & $0$ & $\frac{1}{3}$ & $3$ & & $47$ \\
\hline
$f(x)$ & $20$ & $6$ & & & $-2$ & $-13$ & \\ 
\hline 
\end{tabular} 
\end{center}
\begin{enumerate}
\item Quel est le sens de variation de $f$ ?
\item Déterminer l'expression $f(x)$, puis compléter le tableau.
\end{enumerate}

\exo

\begin{minipage}[t]{9.2cm}
Sur le graphique ci-contre, les droites $\cD_f$, $\cD_g$ et $\cD_h$ représentent des fonctions affines $f$, $g$ et $h$.

Après lecture graphique, construire le tableau de signe de chaque fonction.
\end{minipage}
\hfill
\begin{minipage}[t]{7.3cm}
~

\vspace{-12pt}
\def\xmin{-6}
\def\xmax{6}
\def\ymin{-3}
\def\ymax{4}
\def\xunite{0.6cm}
\def\yunite{0.6cm}
\psset{unit=\xunite, xunit=\xunite, yunit=\yunite, algebraic=true, dimen=middle, dotstyle=*, dotsize=3pt 0, linewidth=0.8pt, arrowsize=3pt 2, arrowinset=0.25, comma=true}
\begin{pspicture*}(\xmin,\ymin)(\xmax,\ymax)
\psgrid[xunit=\xunite, yunit=\yunite, subgriddiv=0, gridlabels=0, gridcolor=lightgray](0,0)(\xmin,\ymin)(\xmax,\ymax)
\psaxes[labels=none,labelFontSize=\scriptstyle, xAxis=true, yAxis=true, Dx=1, Dy=1, ticksize=-2pt 0, subticks=1]{->}(0,0)(\xmin,\ymin)(\xmax,\ymax)[{\footnotesize $x$},135][{\footnotesize $y$},-45]
\psplot[linewidth=1pt, plotpoints=200]{\xmin}{\xmax}{2*x/3+2}
\psplot[linewidth=1pt, plotpoints=200]{\xmin}{\xmax}{-4*x/3+8/3}
\psplot[linewidth=1pt, plotpoints=200]{\xmin}{\xmax}{-1.5}
\psdots(-3,0)(2,0)
\begin{footnotesize}
\uput{3pt}[325](2.25,3.5){$\cD_f$}
\uput{3pt}[215](-0.6,3.5){$\cD_g$}
\uput{3pt}[270](-1.5,-1.5){$\cD_h$}
\uput{3pt}[135](-3,0){$A$}
\uput{3pt}[045](2,0){$B$}
\uput{5pt}[225](0,0){$0$}
\uput{4pt}[270](1,0){$1$}
\uput{4pt}[180](0,1){$1$}
\end{footnotesize}
\end{pspicture*}
\end{minipage}

\exo

Pour chacune des fonctions définies ci-dessous :
\begin{itemize}
\item construire son tableau de signe ;
\item donner un réel $x_1$ dont l'image est positive, et un réel $x_2$ dont l'image est négative.
\end{itemize}

\vspace{-8pt}
\setlength{\columnsep}{1cm}
\setlength{\columnseprule}{0pt}
\begin{multicols}{3}
\begin{enumerate}
\item $f(x)=2x+3\vphantom{\dfrac{h}{p}}$
\item $f(x)=3x-6\vphantom{\dfrac{h}{p}}$
\item $f(x)=-4x+5\vphantom{\dfrac{h}{p}}$
\item $f(x)=x+7\vphantom{\dfrac{h}{p}}$
\item $f(x)=4-5x\vphantom{\dfrac{h}{p}}$
\item $f(x)=8-x\vphantom{\dfrac{h}{p}}$
\item $f(x)=2x-3\vphantom{\dfrac{h}{p}}$
\item $f(x)=\sqrt{3}\;(x-2)\vphantom{\dfrac{h}{p}}$
\item $f(x)=\dfrac{3-2x}{7}\vphantom{\dfrac{h}{p}}$
\item $f(x)=-\dfrac{4}{3}x+8\vphantom{\dfrac{h}{p}}$
\item $f(x)=-2x+\dfrac{1}{2}\vphantom{\dfrac{h}{p}}$
\item $f(x)=\dfrac{x+3}{4}\vphantom{\dfrac{h}{p}}$

\end{enumerate}
\end{multicols}
\vspace{-6pt}

\exo

Résoudre dans $\R$ les inéquations suivantes :

\vspace{-8pt}
\setlength{\columnsep}{1cm}
\setlength{\columnseprule}{0pt}
\begin{multicols}{3}
\begin{enumerate}
\item $2x-2\geqslant 0$
\item $3x-5<0$
\item $4x-3>0$
\end{enumerate}
\end{multicols}

\exo

Une entreprise vend des coques de protection de téléphone.

Le résultat comptable de l'entreprise (bénéfice ou déficit) est donné, en euros, par $R(x)=9x-\np{2583}$, où $x$ est le nombre de coques vendues.
\begin{enumerate}
\item Calculer $R(0)$, et en donner l'interprétation.
\item Construire le tableau de signe de la fonction $R$.
\item Quel est le nombre minimal de coques que l'entreprise doit vendre pour être bénéficiaire ?
\end{enumerate}

\exo

Un congélateur est débranché : sa température intérieure est \SI[retain-explicit-plus]{20}{\celsius}. À $14$ heures, on le branche et sa température diminue alors de \SI[retain-explicit-plus]{1}{\celsius} toutes les $10$ minutes.
\begin{enumerate}
\item Exprimer la température $f(t)$ de l'intérieur du congélateur, où $t$ est la durée (en minutes) écoulée à partir de l'instant du branchement.
\item Construire le tableau de signe de la fonction $f$.
\item À quelle heure la température devient-elle négative ?
\end{enumerate}

\exo

Résoudre dans $\R$ les inéquations suivantes :

\vspace{-8pt}
\setlength{\columnsep}{1cm}
\setlength{\columnseprule}{0pt}
\begin{multicols}{3}
\begin{enumerate}
\item $3x-5\leqslant 7\vphantom{\dfrac{h}{p}}$
\smallskip
\item $2x+5\geqslant -3\vphantom{\dfrac{h}{p}}$
\smallskip
\item $\dfrac{1}{2}x-2\leqslant 6\vphantom{\dfrac{h}{p}}$
\smallskip
\item $-\dfrac{1}{3}x+5\geqslant -4\vphantom{\dfrac{h}{p}}$
\smallskip
\item $\dfrac{4}{3}x-1<\dfrac{1}{3}\vphantom{\dfrac{h}{p}}$
\smallskip
\item $-\dfrac{7}{5}x+2>3\vphantom{\dfrac{h}{p}}$
\smallskip
\item $\dfrac{1}{2}x+4>5\vphantom{\dfrac{h}{p}}$
\smallskip
\item $-\dfrac{1}{5}x+\dfrac{1}{5}\leqslant 0\vphantom{\dfrac{h}{p}}$
\smallskip
\item $-\dfrac{1}{3}x+5\geqslant -4\vphantom{\dfrac{h}{p}}$
\end{enumerate}
\end{multicols}

\exo

Résoudre dans $\R$ les inéquations suivantes :

\vspace{-8pt}
\setlength{\columnsep}{1cm}
\setlength{\columnseprule}{0pt}
\begin{multicols}{3}
\begin{enumerate}
\item $3x+2\leqslant x-14\vphantom{\dfrac{h}{p}}$
\smallskip
\item $-2x-5>4x+31\vphantom{\dfrac{h}{p}}$
\smallskip
\item $-x+19\leqslant -x+51\vphantom{\dfrac{h}{p}}$
\smallskip
\item $-3x+5<-x+17\vphantom{\dfrac{h}{p}}$
\smallskip
\item $4x+\dfrac{1}{3}>2x-\dfrac{1}{3}\vphantom{\dfrac{h}{p}}$
\smallskip
\item $\dfrac{7}{4}x+2>\dfrac{9}{4}x-1\vphantom{\dfrac{h}{p}}$
\end{enumerate}
\end{multicols}

\exo

Résoudre dans $\R$ les inéquations suivantes :

\vspace{-8pt}
\setlength{\columnsep}{1cm}
\setlength{\columnseprule}{0pt}
\begin{multicols}{3}
\begin{enumerate}
\item $2(x+1)-7x>5-x$
\item $4x+5\leqslant 3(x-1)+3$
\item $3(x+4)>x-7$
\end{enumerate}
\end{multicols}

\exo

Un musée propose deux tarifs :
\begin{itemize}
\item Tarif \textbf{A} : chaque entrée coûte $6$ \euro{} ;
\item Tarif \textbf{B} : après avoir payé un abonnement de $16$ \euro{}, chaque entrée coûte $4$ \euro{}.
\end{itemize}
\begin{enumerate}
\item Alice pense aller cinq fois au musée cette année : quel tarif a-t-elle intérêt à choisir ?
\item
\begin{enumerate}
\item Donner les expressions $A(x)$ et $B(x)$ qui indiquent respectivement le prix total à payer pour $x$ entrées, selon que l'on choisit le tarif \textbf{A} ou le tarif \textbf{B}.
\item Résoudre l'inéquation $A(x)>B(x)$.
\item Quelle interprétation peut-on donner du résultat obtenu à la question précédente ?
\end{enumerate}
\end{enumerate}

\exo

Bob a acheté des graines de tomates à $2,90$ \euro{} pour les semer dans son potager.

Il pense revendre des tomates à ses amis au prix de $1,50$ \euro{} le kilogramme, et cherche connaître la quantité $x$ de tomates à revendre pour faire un bénéfice d'au moins $25$ \euro{}.
\begin{enumerate}
\item Modéliser le problème par une inéquation.
\item Résoudre cette inéquation, et répondre au problème de Bob.
\end{enumerate}

\exo

Alice a gagné un concours télévisé. Elle a le choix entre deux lots :
\begin{itemize}
\item une somme de \np{100000} euros{}, puis \np{1400} \euro{} par mois, à vie.
\item une somme de \np{5000} euros{}, puis \np{2000} \euro{} par mois, à vie.
\end{itemize}
Après combien de mois la deuxième offre devient-elle plus intéressante que la première ?

\exo

\begin{minipage}[t]{9.2cm}
Sur la figure ci-contre, les longueurs sont exprimées en centimètres.

Déterminer les valeurs de $x$ pour lesquelles l'aire du domaine coloré dépasse \SI[retain-explicit-plus]{50}{\cm^2}.
\end{minipage}
\hfill
\begin{minipage}[t]{7.3cm}
~

\vspace{-12pt}
\def\xmin{-1}
\def\xmax{9}
\def\ymin{-2.1}
\def\ymax{4.1}
\def\xunite{0.5cm}
\def\yunite{0.5cm}
\psset{unit=\xunite, xunit=\xunite, yunit=\yunite, algebraic=true, dimen=middle, dotstyle=*, dotsize=3pt 0, linewidth=0.8pt, arrowsize=3pt 2, arrowinset=0.25, comma=true}
\begin{pspicture*}(\xmin,\ymin)(\xmax,\ymax)
%\psgrid[xunit=\xunite, yunit=\yunite, subgriddiv=0, gridlabels=0, gridcolor=lightgray](0,0)(\xmin,\ymin)(\xmax,\ymax)
%\psaxes[labels=none,labelFontSize=\scriptstyle, xAxis=true, yAxis=true, Dx=1, Dy=1, ticksize=-2pt 0, subticks=1]{->}(0,0)(\xmin,\ymin)(\xmax,\ymax)[{\footnotesize $x$},135][{\footnotesize $y$},-45]
\pspolygon[fillcolor=lightgray!50, fillstyle=solid](0,0)(2,0)(2,4)(0,4)
\pspolygon[fillcolor=lightgray!50, fillstyle=solid](0,0)(2,0)(2,-2)(0,-2)
\pspolygon[fillcolor=lightgray!50, fillstyle=solid](2,0)(8,2)(2,4)
\psline(2,2)(8,2)
\begin{footnotesize}
\uput{5pt}[180](0,2){$x$}
\uput{5pt}[180](0,-1){$2$}
\uput{5pt}[090](4,2){$6$}
\uput{5pt}[090](1,0){$2$}
\end{footnotesize}
\end{pspicture*}
\end{minipage}

\exo

Il existe différentes unités permettant de mesurer les températures : le degré Celsius (\si{\celsius}), le degré\\ Fahrenheit (\si{\degree}F) utilisé en particulier aux États-Unis, le Kelvin (\si{\kelvin}) utilisé par les physiciens.

Le tableau suivant donne quelques équivalences entre ces unités :
\begin{center}
\setlength{\tabcolsep}{0.1cm}
\renewcommand{\arraystretch}{1.2}
\begin{tabular}{
>{\centering\arraybackslash}m{5cm}
*{3}{>{\centering\arraybackslash}m{2cm}}}
\hline
zéro absolu & \SI{-273}{\celsius} &  & \SI{0}{\kelvin} \\ 
\hline 
fusion de l'eau & \SI{0}{\celsius} & \SI{32}{\degree}F & \SI{273}{\kelvin} \\ 
\hline 
vaporisation de l'eau & \SI{100}{\celsius} & \SI{212}{\degree}F & \\ 
\hline
\end{tabular} 
\end{center}
D'autre part, ces trois échelles sont liées par des relations affines.
\begin{enumerate}
\item Si $x$ est une température exprimée en \si{\celsius}, quelle est sa valeur $a(x)$ en \si{\degree}F ?
\item Si $x$ est une température exprimée en \si{\degree}F, quelle est sa valeur $b(x)$ en \si{\celsius} ?
\item Si $x$ est une température exprimée en \si{\degree}F, quelle est sa valeur $c(x)$ en \si{\kelvin} ?
\end{enumerate}

\exo

\begin{enumerate}
\item Tracer dans un repère la représentation graphique de la fonction $f$ définie sur $\R$ par :
\[f(x)=
\begin{cases}
2x+5 & \textrm{si }x\leqslant -2 \\
-x-1 & \textrm{si }-2<x\leqslant 4 \\
2x-13 & \textrm{si }x>4 
\end{cases}\]
\item Écrire une fonction Python qui prend en argument un réel $x$, et renvoie son image par la fonction $f$.
\end{enumerate}

